\section*{MEASURES ON GENERAL TOPOLOGICAL SPACES AND STRICT CONVERGENCE.}
\phantomsection
\addcontentsline{toc}{section}{MEASURES ON GENERAL TOPOLOGICAL SPACES AND STRICT CONVERGENCE.}

The study of the links between topology and measure theory was especially thought of as the study of regularity properties of measures, and in particular that of ``exterior'' regularity and of ``interior'' regularity; \(^{14}\) interior regularity is equivalent to exterior regularity on a locally compact space countable at infinity. The construction that Lebesgue gives of the measure of sets on the straight line brings out these two kinds of regularity, and the property of exterior regularity of measures on a Polish space seems to have had public notoriety around 1935. But it is only in 1940, in an article whose diffusion was slowed down by the war, that A. D. Alexandroff {[}3{]} brought out the role of interior regularity and showed that this is possessed by the measures on a Polish space; this result was found again later by Prokhorov {[}254{]} and is often attributed mistakenly to this author. It only appeared very recently that this property extended to sublinian spaces; because of this, the importance of these spaces increased greatly, all the more so that it became known that their theory could be done without a metrisability hypothesis, and that the quasi-totality of functional spaces were sublinian (and even most often lusinian). \(^{15}\)

The definition of a mode of convergence (weak or strict) for measures is done in the most convenient way by putting in duality the space of measures with a space of continuous functions. Generalising an ancient result of F. Riesz, A. A. Markhoff established in 1938 a bijective correspondence between the positive functionals on \(\mathcal{C}(X)\) and the regular measures on a compact space X. In the work {[}3{]} already quoted, A. D.. Alexandroff extends these results to the case of a completely regular space: he introduces a hierarchy on the set of positive linear forms on the space \(\mathcal{C}^{b}(X)\) of bounded continuous functions on a completely regular space X, \(^{16}\) he defines strict convergence of bounded measures and proves among others the following two theorems:

\begin{enumerate}
\def\labelenumi{\alph{enumi})}
\item
  if X is Polish, the set of linear forms on \(\mathcal{C}^{b}(X)\) corresponding to measures is closed in the weak convergence of sequences;
\item
  if a sequence of bounded measures has a strict limit, ``there is no mass fleeing to infinity'' (it is a weak form of the reciprocal of the theorem of strict convergence of Prokhorov).
\end{enumerate}

From this fusion of notions and theorems, Prokhorov will know how to extract important results for the theory of stochastic processes, and to present them in a simple and striking form. In his great work of 1956 already quoted {[}254{]}, an important part is devoted to bounded positive measures on a Polish space; in generalising a construction of Lévy, he defines on the set of positive measures of mass 1 a distance which makes it into a Polish space, then he establishes an important compactness criterion for strict convergence. Independently of Prokhorov, Le Cam {[}197{]} has obtained a certain number of compactness results for the strict convergence of measures; he puts no metrisability hypotheses on the spaces that he considers, and his results reduce to previous theorems of Dieudonné in the locally compact case.

